\section{Layer~2 Results}\label{sec:l2res}
\subsection{Reference versus compiled PQC}
Table~\ref{tab:pqc-comparison} compares median latencies. Over all compared operations the speedup ranged from \SpeedupMin{}$\times$ to \SpeedupMax{}$\times$. ML-KEM-768 key generation, encapsulation and decapsulation were \SpeedupMlkemKeygen{}$\times$, \SpeedupMlkemEncaps{}$\times$ and \SpeedupMlkemDecaps{}$\times$ faster with the compiled backend, and ML-DSA-65 key generation and verification were \SpeedupMldsaKeygen{}$\times$ and \SpeedupMldsaVerify{}$\times$ faster, i.e.\ these costs are strongly implementation-dependent. ML-DSA-65 signing improved only \SpeedupMldsaSign{}$\times$ (median \RefMldsaSignMs{}~ms to \OptMldsaSignMs{}~ms), consistent with the variable-iteration rejection sampling in its signing procedure.

\begin{table}[t]
\centering
\small
\caption{Reference Python vs.\ \texttt{pqcrypto}: median latency and implementation-backend speedup (selected parameter sets). $n$ is the number of timed repetitions.}
\label{tab:pqc-comparison}
\resizebox{\linewidth}{!}{%
\begin{tabular}{llrrrrr}
\toprule
Algorithm & Op. & n ref & n opt & Ref. median (ms) & pqcrypto median (ms) & Speedup \\
\midrule
ML-DSA-65 & keygen & 30 & 200 & 11.5 & 0.237 & 48.5 \\
ML-DSA-65 & sign & 30 & 200 & 45.5 & 9.45 & 4.81 \\
ML-DSA-65 & verify & 30 & 200 & 12.7 & 0.22 & 57.6 \\
ML-KEM-768 & decaps & 30 & 200 & 6.78 & 0.111 & 61 \\
ML-KEM-768 & encaps & 30 & 200 & 5.02 & 0.0876 & 57.3 \\
ML-KEM-768 & keygen & 30 & 200 & 3.71 & 0.0721 & 51.5 \\
SLH-DSA-SHA2-128f & keygen & 3 & 200 & 22.8 & 2.56 & 8.94 \\
SLH-DSA-SHA2-128f & sign & 3 & 200 & 231 & 60.2 & 3.84 \\
SLH-DSA-SHA2-128f & verify & 10 & 200 & 12.8 & 3.68 & 3.47 \\
SLH-DSA-SHA2-128s & keygen & 3 & 30 & 1.25e+03 & 163 & 7.69 \\
SLH-DSA-SHA2-128s & sign & 3 & 30 & 4.33e+03 & 1.27e+03 & 3.41 \\
SLH-DSA-SHA2-128s & verify & 10 & 200 & 4.42 & 1.23 & 3.59 \\
\bottomrule
\end{tabular}}
\end{table}


\paragraph{The SLH-DSA question.} An earlier session measured about 15~s per SLH-DSA-SHA2-128s signature with the reference code (n=5); that value is preserved. In the controlled run reported here, the reference median was \RefSlhSSignMs{}~ms (n=\SlhSRefN{}) and the compiled median \OptSlhSSignMs{}~ms (n=\SlhSOptN{}), a \SpeedupSlhSSign{}$\times$ difference, and the fast variant 128f signed in \OptSlhFSignMs{}~ms compiled. Two conclusions follow. First, the ${\sim}15$~s figure is reference-implementation behaviour from a different session and is not an implementation-independent property. Second, compiled SLH-DSA-128s signing remains about two orders of magnitude slower than compiled ML-DSA-65 (\OptMldsaSignMs{}~ms), so most of its signing cost is a property of the parameter set rather than only of the implementation.

\begin{figure}[t]
\centering
\includegraphics[width=\figw]{pqc_backend_speedup}
\caption{Implementation-backend speedup per operation (median ratio). This compares implementations, not cryptographic strength.}
\label{fig:speedup}
\end{figure}
\begin{figure}[t]
\centering
\includegraphics[width=\figw]{pqc_ref_vs_opt_slhdsa}
\caption{SLH-DSA reference versus compiled latency (log-scale; reference signing $n=3$).}
\label{fig:slh}
\end{figure}
\begin{figure}[t]
\centering
\includegraphics[width=\figw]{classical_vs_pqc_latency}
\caption{Median latency by cryptographic role, classical versus post-quantum (log-scale). AES-256 is excluded because it performs a different function.}
\label{fig:classical}
\end{figure}
\begin{figure}[t]
\centering
\includegraphics[width=\figw]{crypto_object_sizes}
\caption{Key, ciphertext and signature sizes (log-scale; missing bar means not applicable). Sizes are identical across the two PQC backends.}
\label{fig:sizes}
\end{figure}

\subsection{End-to-end banking-transaction protection}
Table~\ref{tab:e2e} gives total latency including serialization, signing, key establishment, encryption, decryption and verification. The classical configuration had a median of \EEAMedian{}~ms; compiled ML-KEM-768 + ML-DSA-65 \EECMedian{}~ms; the same with reference code \EEBMedian{}~ms; the hybrid configuration \EEDMedian{}~ms. Configurations using SLH-DSA-128s were dominated by signing (\EEBTwoMedian{}~ms reference, \EECTwoMedian{}~ms compiled). Secured payload grew to \EEAExpansion{}$\times$ the plaintext (\PlaintextBytes{}~B) classically, \EECExpansion{}$\times$ with ML-KEM/ML-DSA and \EECTwoExpansion{}$\times$ with ML-KEM/SLH-DSA, almost entirely from KEM ciphertext and signature bytes. Simulated transmission is in-memory, so no network latency is included.

\begin{table}[t]
\centering
\small
\caption{End-to-end banking-transaction protection (serialize--sign--key establishment--AES-256-GCM--decrypt--verify). Plaintext is the same canonical JSON payload in all rows.}
\label{tab:e2e}
\resizebox{\linewidth}{!}{%
\begin{tabular}{lrrrrr}
\toprule
Configuration & n & Median (ms) & P95 (ms) & Secured (B) & Expansion \\
\midrule
A\_classical\_ECDH+ECDSA & 200 & 0.302 & 0.568 & 475 & 1.52 \\
B\_reference\_ML-KEM-768+ML-DSA-65 & 30 & 72.8 & 141 & 4737 & 15.2 \\
C\_optimized\_ML-KEM-768+ML-DSA-65 & 200 & 10.2 & 11 & 4737 & 15.2 \\
D\_hybrid\_optimized\_ECDH+ML-KEM-768\_ECDSA+ML-DSA-65 & 200 & 10.5 & 11 & 4873 & 15.6 \\
B2\_reference\_ML-KEM-768+SLH-DSA-128s & 3 & 3.63e+03 & 3.66e+03 & 9284 & 29.8 \\
C2\_optimized\_ML-KEM-768+SLH-DSA-128s & 20 & 1.25e+03 & 1.27e+03 & 9284 & 29.8 \\
\bottomrule
\end{tabular}}
\end{table}


\begin{figure}[t]
\centering
\includegraphics[width=\figw]{end_to_end_latency}
\caption{End-to-end latency by configuration (log-scale; whiskers to P95).}
\label{fig:e2e}
\end{figure}

\subsection{BB84 statistics and cross-validation}
Table~\ref{tab:bb84} and Fig.~\ref{fig:qberci} report QBER with confidence intervals. The clean channel gave QBER 0 in both implementations. With injected noise of 1\%, 5\% and 10\%, mean QBER was \BbqProtoKOnePct{}, \BbqProtoKFivePct{} and \BbqProtoKTenPct{} (protocol, 2000 bits) and \BbqQiskOnePct{}, \BbqQiskFivePct{} and \BbqQiskTenPct{} (Qiskit, 500 bits). With the intercept--resend attacker, mean QBER was \BbqProtoKEve{} (95\% CI \BbciProtoKEve{}) and \BbqQiskEve{} (95\% CI \BbciQiskEve{}) against the theoretical 0.25; adding 2\% noise raised them to \BbqProtoKEveNoise{} and \BbqQiskEveNoise{}. Both attack intervals can be read directly against 0.25 in Fig.~\ref{fig:qberci}; no value was tuned toward the theory. Comparing the two implementations at 500 bits per trial gave MAE \BbMAE{}, RMSE \BbRMSE{} and a maximum difference of \BbMaxDiff{} in mean QBER. Sifting efficiency was \SiftEff{} (theory 0.5).

\begin{table}[t]
\centering
\small
\caption{BB84 software simulation: protocol simulator (2000 bits/trial) and Qiskit/Aer circuit simulator (500 bits/trial). QBER over the full sifted key; 95\% $t$-interval over trials; detection rate = fraction of trials with QBER above 0.11.}
\label{tab:bb84}
\resizebox{\linewidth}{!}{%
\begin{tabular}{llrrrrrrr}
\toprule
Impl. & Condition & Trials & Bits & \$\textbackslash eta\_\{sift\}\$ & Mean QBER & CI low & CI high & Det. rate \\
\midrule
protocol 2000 & clean & 100 & 2000 & 0.4999 & 0.0000 & 0.0000 & 0.0000 & 0.0000 \\
protocol 2000 & noise\_1pct & 100 & 2000 & 0.4999 & 0.0097 & 0.0091 & 0.0103 & 0.0000 \\
protocol 2000 & noise\_5pct & 100 & 2000 & 0.4999 & 0.0490 & 0.0477 & 0.0504 & 0.0000 \\
protocol 2000 & noise\_10pct & 100 & 2000 & 0.4999 & 0.0979 & 0.0960 & 0.0999 & 0.1200 \\
protocol 2000 & eve\_intercept\_resend & 100 & 2000 & 0.4996 & 0.2508 & 0.2479 & 0.2537 & 1.0000 \\
protocol 2000 & eve\_and\_noise\_2pct & 100 & 2000 & 0.4996 & 0.2606 & 0.2578 & 0.2635 & 1.0000 \\
qiskit 500 & clean & 50 & 500 & 0.4991 & 0.0000 & 0.0000 & 0.0000 & 0.0000 \\
qiskit 500 & noise\_1pct & 50 & 500 & 0.4991 & 0.0082 & 0.0059 & 0.0104 & 0.0000 \\
qiskit 500 & noise\_5pct & 50 & 500 & 0.4991 & 0.0498 & 0.0458 & 0.0538 & 0.0000 \\
qiskit 500 & noise\_10pct & 50 & 500 & 0.4991 & 0.0960 & 0.0896 & 0.1023 & 0.2200 \\
qiskit 500 & eve\_intercept\_resend & 50 & 500 & 0.4991 & 0.2438 & 0.2360 & 0.2515 & 1.0000 \\
qiskit 500 & eve\_and\_noise\_2pct & 50 & 500 & 0.4991 & 0.2649 & 0.2576 & 0.2722 & 1.0000 \\
\bottomrule
\end{tabular}}
\end{table}


\begin{figure}[t]
\centering
\includegraphics[width=\figw]{bb84_final_qber_ci}
\caption{Mean QBER with 95\% confidence intervals by condition and implementation; dotted line: 25\% theory for intercept--resend.}
\label{fig:qberci}
\end{figure}
\begin{figure}[t]
\centering
\includegraphics[width=\figw]{bb84_final_eve_distribution}
\caption{Per-trial QBER with and without intercept--resend; dashed line: 0.11 threshold.}
\label{fig:evedist}
\end{figure}

\subsection{Eve-detection threshold}
At the conventional 0.11 threshold, the true-positive rate was \ThrTPRProtoK{} (protocol, 2000 bits), \ThrTPRProtoH{} (protocol, 500 bits) and \ThrTPRQisk{} (Qiskit), while the false-positive rate when 10\%-noise trials count as Eve-absent was \ThrFPRProtoK{}, \ThrFPRProtoH{} and \ThrFPRQisk{} (Table~\ref{tab:threshold}, Fig.~\ref{fig:thr}). A higher threshold (first maximum of balanced accuracy at \ThrBestAProtoK{}, \ThrBestAProtoH{} and \ThrBestAQisk{}) separates the classes in these simulations, but that optimum depends on the noise range chosen and is not a general security threshold. QBER above a threshold does not prove eavesdropping, because channel noise also increases QBER.

\begin{table}[t]
\centering
\small
\caption{Eve-detection threshold analysis at the conventional 0.11 value; last column is the lowest threshold attaining the maximum balanced accuracy in this simulated design.}
\label{tab:threshold}
\resizebox{\linewidth}{!}{%
\begin{tabular}{llrrrr}
\toprule
Impl. & Eve-absent set & TPR@0.11 & FPR@0.11 & Miss@0.11 & First best thr. \\
\midrule
protocol 2000 & no eve all noise levels & 1.000 & 0.030 & 0.000 & 0.120 \\
protocol 2000 & no eve clean to 5pct & 1.000 & 0.000 & 0.000 & 0.070 \\
protocol 500 & no eve all noise levels & 1.000 & 0.085 & 0.000 & 0.140 \\
protocol 500 & no eve clean to 5pct & 1.000 & 0.000 & 0.000 & 0.080 \\
qiskit 500 & no eve all noise levels & 1.000 & 0.055 & 0.000 & 0.160 \\
qiskit 500 & no eve clean to 5pct & 1.000 & 0.000 & 0.000 & 0.090 \\
\bottomrule
\end{tabular}}
\end{table}


\begin{figure}[t]
\centering
\includegraphics[width=\figw]{bb84_threshold_tradeoff}
\caption{Eve-detection threshold trade-off: TPR/FPR against threshold and ROC-style view.}
\label{fig:thr}
\end{figure}
